\chapter{METODOLOGI}

Bab ini menjelaskan protokol penelitian yang diusulkan untuk mengukur perubahan kinerja klasifikasi arah XAUUSD setelah penambahan kelompok fitur teknikal dan makroekonomi. Seluruh ambang, periode, fitur, dan hiperparameter ditetapkan sebagai pilihan desain sebelum eksperimen utama, bukan sebagai nilai optimum yang telah dibuktikan. Arsip uji kelayakan yang tersedia menjadi bukti bahwa penelitian dapat dilanjutkan, tetapi hasilnya tidak diperlakukan sebagai hasil eksperimen utama ataupun sebagai data uji prospektif yang belum pernah dilihat.

\section{Tipe Penelitian}

Penelitian ini bertipe nonimplementatif--analitik dalam bidang komputasi cerdas. Tugas klasifikasi digunakan sebagai instrumen untuk membandingkan kandungan informasi tiga konfigurasi fitur, bukan untuk membangun aplikasi transaksi, memberikan rekomendasi perdagangan, atau menjelaskan hubungan sebab-akibat ekonomi. Penelitian tidak melibatkan partisipan manusia.

\section{Strategi Penelitian}

Strategi penelitian berupa eksperimen komparatif dengan penambahan kelompok fitur secara kumulatif. Dua model dengan karakter berbeda digunakan, yaitu Regresi Logistik Multinomial sebagai pembanding linear dan XGBoost sebagai pembanding nonlinear berbasis pohon. Pemakaian kedua model tersebut mengikuti preseden klasifikasi arah aset dengan model linear dan berbasis pohon \parencite{sadorsky2021predicting,basher2022forecasting}. Kedua model diterapkan pada K1, K2, dan K3 untuk menilai apakah perubahan kinerja antarkonfigurasi konsisten pada karakter model yang berbeda. Perbedaan skor kedua model dilaporkan secara deskriptif, bukan sebagai uji linearitas data atau dasar untuk menyatakan satu model selalu lebih unggul.

Perubahan skor Konfigurasi 2 terhadap Konfigurasi 1 ($K2-K1$) dan Konfigurasi 3 terhadap Konfigurasi 2 ($K3-K2$) dihitung secara berpasangan pada bar uji, periode uji, model, dan metrik yang sama. Pada analisis utama, setiap konfigurasi memperoleh penyetelan hiperparameter sendiri menggunakan data latih. Dengan demikian, selisih skor mengukur perubahan bersyarat akibat penambahan kelompok fitur beserta penyetelan ulang model, bukan pengaruh murni atau kausal fitur tersebut. Analisis sensitivitas dengan hiperparameter tetap digunakan untuk membantu memisahkan kedua unsur itu.

Alur penelitian yang direncanakan ditunjukkan pada Gambar \ref{fig:alur-penelitian}. Protokol dibekukan sebelum eksperimen baru dijalankan dan sebelum skor periode uji diperiksa.

\begin{figure}[htbp]
    \centering
    \begin{tikzpicture}[
        node distance=3.5mm,
        every node/.style={font=\small,align=center},
        process/.style={rectangle,draw=black,rounded corners=1.5pt,
            minimum width=10.8cm,minimum height=0.83cm,text width=10.2cm,
            inner xsep=4pt,inner ysep=3pt},
        flow/.style={-{Latex[length=2.2mm]},thin}
    ]
        \node[process] (input) {Mulai: data harga dan kalender ekonomi sekunder};
        \node[process,below=of input] (audit) {Audit kualitas harga dan kalender\\serta verifikasi integritas dan waktu rilis};
        \node[process,below=of audit] (feature) {Pelabelan kronologis dan pembentukan fitur\\harga, teknikal, dan makroekonomi};
        \node[process,below=of feature] (inner) {Validasi internal: praproses, penyetelan,\\dan pelatihan ulang pada data latih luar};
        \node[process,below=of inner] (outer) {Prediksi pada periode uji utama dan penghitungan\\macro-F1, balanced accuracy, serta matriks kebingungan};
        \node[process,below=of outer] (compare) {Perbandingan berpasangan $K2-K1$ dan $K3-K2$\\serta analisis sensitivitas};
        \draw[flow] (input) -- (audit);
        \draw[flow] (audit) -- (feature);
        \draw[flow] (feature) -- (inner);
        \draw[flow] (inner) -- (outer);
        \draw[flow] (outer) -- (compare);
    \end{tikzpicture}
    \caption{Alur penelitian yang diusulkan}
    \label{fig:alur-penelitian}
\end{figure}

\section{Sumber dan Teknik Pengumpulan Data}

\subsection{Sumber dan Cakupan Data}

Penelitian menggunakan data sekunder yang direkam oleh vendor. Data harga utama berupa bar \emph{open}, \emph{high}, \emph{low}, dan \emph{close} (OHLC) XAUUSD H1 dari server MetaQuotes-Demo. Arsip kelayakan yang telah diverifikasi mencakup 38.799 bar dari Januari 2020 sampai Juli 2026 atau sekitar 6,578 tahun. Waktu bar yang dikembalikan antarmuka Python MetaTrader 5 diperlakukan sebagai UTC sesuai dokumentasi resminya.\footnote{Dokumentasi resmi \emph{copy rates range}: \url{https://www.mql5.com/en/docs/python_metatrader5/mt5copyratesrange_py}.} Jumlah observasi eksperimen akhir ditentukan setelah seluruh aturan kelayakan, jeda, pemanasan indikator, dan pelabelan diterapkan.

Data makroekonomi berasal dari kalender ekonomi internal MetaTrader 5. Snapshot arsip beku ini berisi 706 kejadian USD yang terverifikasi tanpa duplikasi pengenal nilai, dan keluarga kejadian seperti NFP dan tingkat pengangguran memiliki data aktual dan prakiraan utuh yang secara struktural menutupi periode data latih dan setiap periode uji utama. Dokumentasi kalender menyatakan bahwa waktu kalender menggunakan waktu server perdagangan.\footnote{Dokumentasi resmi \emph{CalendarValueHistory}: \url{https://www.mql5.com/en/docs/calendar/calendarvaluehistory}.} Namun, apabila penanda waktu mentah historis pada snapshot beku spesifik ini dilabeli ulang sebagai UTC, seluruh sepuluh jangkar waktu rilis resmi telah terverifikasi cocok secara tepat; oleh sebab itu, rentang waktu arsip kalender ini diperlakukan sebagai UTC dan dikontrol ulang menggunakan pemeriksa hash SHA-256 dan penanda waktu pada Subbab \ref{subsec:audit-kalender}.

Keluarga peristiwa dibatasi pada CPI umum bulanan, CPI umum tahunan, CPI inti bulanan, CPI inti tahunan, \emph{Nonfarm Payrolls} (NFP), tingkat pengangguran, keputusan suku bunga, dan peristiwa FOMC. Pemilihan ini didukung oleh temuan bahwa kejutan makroekonomi berkaitan dengan pembentukan harga emas intrahari \parencite{sobti2021how}.

\subsection{Teknik Pengumpulan Data}

Bar harga diperoleh secara baca-saja melalui antarmuka Python MetaTrader 5. Ekstraksi arsip awal dilakukan pada Windows 11 tanpa pemanggilan API pesanan atau perdagangan. Hasil ekstraksi disimpan sebagai CSV dan Parquet bersama manifes yang mencatat sumber, simbol, rentang waktu, serta checksum berkas. Windows dan terminal MT5 hanya diperlukan kembali apabila ekstraktor harga Python digunakan untuk mengambil sumber atau rentang harga baru; penyuntingan proposal dan analisis luring atas arsip yang tersedia tidak memerlukannya.

Kalender ekonomi diperoleh secara baca-saja melalui eksportir MQL5 yang telah berfungsi pada MT5 dalam lingkungan Wine di macOS. CSV hasil ekspor mempertahankan kolom waktu mentah, nilai aktual, dan nilai prakiraan agar provenance serta pemeriksaan waktu dapat dilakukan. Pengambilan ini tidak menggunakan kredensial di dalam kode, tidak mengaktifkan transaksi otomatis, dan tidak memanggil API pesanan.

Data harga dan kalender digunakan secara privat dan luring. Berkas mentah dipertahankan terpisah dari data turunan; keluaran pemrosesan tidak menimpa sumber. Sebelum digunakan, identitas sumber, rentang waktu, skema, jumlah bar atau kejadian, checksum, duplikasi, urutan waktu, dan kelengkapan kolom diperiksa serta dicatat dalam manifes audit.

\section{Lingkungan dan Peralatan Penelitian}

Penelitian tidak memerlukan lokasi fisik atau laboratorium tertentu karena penyiapan dan analisis data dilakukan secara luring pada komputer lokal. Komputasi dirancang untuk CPU dan tidak mensyaratkan GPU. Lingkungan persiapan yang tersedia telah diukur pada 12 September 2026 sebagaimana Tabel \ref{tab:perangkat-keras}; pencatatan ini tidak menyatakan bahwa eksperimen utama telah dilaksanakan.

\begin{table}[htbp]
    \centering
    \caption{Perangkat keras lingkungan persiapan penelitian}
    \label{tab:perangkat-keras}
    \begin{tabular}{ll}
        \hline
        Komponen & Spesifikasi terukur \\
        \hline
        Prosesor & Apple M5 (arsitektur arm64) \\
        Memori & 24 GiB RAM \\
        Akselerator & Tidak diperlukan; analisis berbasis CPU \\
        \hline
    \end{tabular}
\end{table}

Perangkat lunak lingkungan persiapan diringkas pada Tabel \ref{tab:perangkat-lunak}. Versi sistem operasi, Python, pustaka, dan Git pada pelaksanaan akhir akan dicatat kembali bersama artefak eksperimen. Lingkungan sumber historis arsip harga tetap dicatat secara terpisah, yaitu Windows 11 dengan Python 3.12.10, NumPy 2.5.1, pandas 3.0.5, scikit-learn 1.9.0, dan XGBoost 2.1.4; pencatatan lingkungan macOS tidak menggantikan provenance tersebut.

\begin{table}[htbp]
    \centering
    \caption{Perangkat lunak lingkungan persiapan penelitian}
    \label{tab:perangkat-lunak}
    \begin{tabular}{ll}
        \hline
        Komponen & Versi terukur pada 12 September 2026 \\
        \hline
        Sistem operasi & macOS 26.6.2 \\
        Bahasa pemrograman & Python 3.14.5 pada lingkungan virtual proyek \\
        Pustaka numerik dan data & NumPy 2.5.1; pandas 3.0.5; PyArrow 23.0.1 \\
        Pustaka pemodelan & scikit-learn 1.9.0; XGBoost 2.1.4 \\
        \hline
    \end{tabular}
\end{table}

\section{Unit Observasi, Waktu Keputusan, dan Target}

Penanda waktu pada data harga menyatakan waktu pembukaan bar. Misalkan $o_t$ adalah waktu pembukaan bar ke-$t$, sedangkan $O_t$, $H_t$, $L_t$, dan $C_t$ berturut-turut adalah harga pembukaan, tertinggi, terendah, dan penutupannya. Keputusan untuk observasi ke-$t$ dibuat setelah bar tersebut selesai, yaitu
\begin{equation}
\tau_t=o_t+1\text{ jam}.
\label{eq:waktu-keputusan}
\end{equation}
Semua fitur harga dan teknikal harus telah tersedia paling lambat pada $\tau_t$. Peristiwa makro hanya digunakan jika waktu rilis terverifikasinya lebih kecil secara ketat daripada $\tau_t$; rilis yang tepat sama dengan $\tau_t$ ditunda ke keputusan berikutnya.

Target merupakan imbal hasil penutupan satu bar berikutnya,
\begin{equation}
q_{t+1}=\ln\left(\frac{C_{t+1}}{C_t}\right).
\label{eq:target-satu-jam}
\end{equation}
Data harga yang diperoleh tidak memuat label BUY, HOLD, atau SELL. Peneliti membentuk label tersebut secara deterministik dari $C_t$ dan $C_{t+1}$: BUY untuk kenaikan di atas ambang positif, HOLD untuk perubahan di antara kedua ambang, dan SELL untuk penurunan di bawah ambang negatif. Ketiganya menyatakan kelas arah, bukan perintah transaksi. Rumus ambang dan syarat kelayakan bar dijelaskan berikut ini.

Observasi hanya layak diberi target apabila bar berikutnya dibuka tepat pada $o_{t+1}=o_t+1$ jam dan target berakhir pada $\tau_{t+1}=o_t+2$ jam. Dengan demikian, label tidak menjembatani akhir pekan, penutupan pasar, atau celah data. Pada setiap periode uji, $\tau_{t+1}$ juga harus berada di dalam interval uji yang sama.

Ambang kelas utama ditetapkan sebesar setengah volatilitas historis 24 bar. Dengan $r_t$ sebagaimana Persamaan \ref{eq:logreturn}, rata-rata dan simpangan baku sampel dihitung sebagai
\begin{align}
\bar r_{t,24} &= \frac{1}{24}\sum_{i=0}^{23}r_{t-i},\\
\sigma_{t,24} &= \sqrt{\frac{1}{23}\sum_{i=0}^{23}\left(r_{t-i}-\bar r_{t,24}\right)^2},\\
\delta_t &= 0{,}5\sigma_{t,24}.
\label{eq:ambang-volatilitas}
\end{align}
Nilai $r_t$ sampai $r_{t-23}$ telah diketahui pada $\tau_t$. Kelas kemudian ditentukan dengan mengganti $r_{t+1}$ pada Persamaan \ref{eq:labeling} menjadi $q_{t+1}$. Observasi dengan $\sigma_{t,24}$ nol, tidak hingga, atau tidak dapat dihitung dikeluarkan. Tidak ada imputasi pada OHLC maupun label. Pengali 0,5 tidak disetel menurut konfigurasi atau periode uji. Pengali 0,25 dan 1,0 diterapkan sebagai dua analisis sensitivitas terpisah pada seluruh konfigurasi, tanpa memilih pengali berdasarkan skor terbaik.

\section{Audit dan Penyiapan Data}

\subsection{Audit Data Harga}

Data harga diperiksa terhadap duplikasi waktu, urutan waktu, nilai yang hilang, harga nonpositif, dan pelanggaran hubungan $H_t\geq\max(O_t,C_t)$ serta $L_t\leq\min(O_t,C_t)$. Audit celah menggunakan jadwal perdagangan yang berlaku pada sumber data. Penutupan pasar yang dapat dijelaskan dicatat sebagai penutupan wajar dan jendela fitur dilanjutkan dengan bar perdagangan sah sebelumnya. Dengan aturan ini, volatilitas fitur merangkum return antarbar perdagangan, termasuk perubahan melintasi penutupan wajar, sedangkan horizon target tetap tepat satu jam.

Celah yang tidak dapat dijelaskan memutus keadaan indikator sehingga perhitungan bergulir, Wilder, dan EMA dimulai kembali setelah celah. Bar OHLC tidak diinterpolasi. Aturan kelayakan ditentukan berdasarkan kualitas data dan ketersediaan informasi, bukan skor model.

\subsection{Audit Data Kalender}
\label{subsec:audit-kalender}

Sebagai kontrol reproduksibilitas untuk snapshot spesifik ini, sebelum data kalender digabungkan, berkas harus cocok dengan hash SHA-256 arsip beku yang ditetapkan pada konfigurasi. Kolom UTC baru dibuat tanpa menimpa waktu mentah. Hasil konversinya dipastikan cocok tepat dengan sepuluh jangkar rilis resmi yang mencakup periode \emph{daylight saving time} dan non-DST. Semantik waktu dan kontrol jangkar rilis historis ini diberlakukan eksklusif untuk snapshot ini dan tidak digeneralisasikan pada broker, \emph{build}, atau snapshot kalender lainnya.

Penelitian ini menggunakan asumsi operasional bahwa kolom \emph{forecast} pada rekaman MT5 merupakan nilai prakiraan yang diasosiasikan vendor dengan kejadian sebelum nilai aktual diketahui. Asumsi ini konsisten dengan dokumentasi kalender MQL5 yang membedakan nilai prakiraan dari nilai aktual yang tersedia setelah publikasi.\footnote{Dokumentasi resmi konsep kalender MQL5: \url{https://www.mql5.com/en/book/advanced/calendar/calendar_overview}.} Riwayat perubahan independen tidak tersedia sehingga snapshot ini tidak diklaim sebagai arsip \emph{point-in-time}. Nilai \emph{previous} tidak digunakan. Untuk setiap keluarga rilis, kejadian terbaru menggantikan keadaan kejadian sebelumnya. Apabila angka pada kolom nilai kejadian kosong, nilai numeriknya tidak direkayasa menjadi nol melainkan dibiarkan tidak tersedia.

\subsection{Prapemrosesan dan Ketertelusuran}

Nilai numerik yang tidak diketahui tidak diubah menjadi nol. Pada setiap blok validasi internal dan periode uji utama, median dihitung hanya dari data latih, kemudian digunakan untuk mengisi nilai hilang pada data latih dan uji bersama indikator biner ketiadaan. Kolom yang seluruhnya hilang pada data latih dibuang dan dicatat. Daftar kolom ditentukan hanya dari data latih dan dipertahankan pada data uji. Ketersediaan angka pada data uji dilaporkan sebagai cakupan evaluasi tanpa digunakan untuk memilih ulang kolom atau menyetel model.

Arsip kelayakan telah diperiksa ketika rancangan disusun. Oleh sebab itu, periode 2020--Juli 2026 tidak diklaim sebagai \emph{holdout} prospektif yang sama sekali belum pernah terlihat. Perlindungan terhadap penyesuaian setelah melihat data dilakukan dengan membekukan definisi pada bab ini sebelum eksperimen baru, menyimpan seluruh kegagalan, dan tidak mengubah protokol berdasarkan skor periode uji.

\section{Pembentukan Fitur}

\subsection{Fitur Harga}

Kelompok harga terdiri atas sembilan fitur: enam kelambanan imbal hasil $r_t,r_{t-1},\ldots,r_{t-5}$, bentuk bar $\ln(C_t/O_t)$, rentang $\ln(H_t/L_t)$, dan $\sigma_{t,24}$. Seluruhnya telah diketahui pada $\tau_t$. Nilai OHLC mentah tidak digunakan sebagai pengganti fitur yang hilang dan tidak diimputasi.

\subsection{Fitur Teknikal}

Kelompok teknikal terdiri atas empat fitur: RSI 14 periode serta garis, sinyal, dan histogram MACD $(12,26,9)$. Pemilihan RSI dan MACD mengikuti penggunaannya sebagai indikator momentum dan tren \parencite{wilder1978new,murphy1999technical}. Tidak ada SMA opsional atau pencarian panjang jendela pada data uji.

Untuk RSI, perubahan penutupan didefinisikan sebagai $d_t=C_t-C_{t-1}$, dengan kenaikan $g_t=\max(d_t,0)$ dan penurunan $\ell_t=\max(-d_t,0)$. Rata-rata kenaikan dan penurunan awal adalah rata-rata aritmetika dari 14 perubahan pertama. Nilai berikutnya diperbarui dengan penghalusan Wilder berparameter $\alpha=1/14$. RSI dihitung sebagai $100-100/(1+\overline g_t/\overline\ell_t)$. Jika kedua rata-rata nol, RSI ditetapkan 50; jika hanya rata-rata penurunan yang nol, RSI ditetapkan 100. Jika hanya rata-rata kenaikan yang nol, rumus menghasilkan RSI 0.

EMA periode $p$ diinisialisasi dengan rata-rata $p$ penutupan pertama dan diperbarui menggunakan $\alpha_p=2/(p+1)$. Garis MACD adalah $\operatorname{EMA}_{12}(C_t)-\operatorname{EMA}_{26}(C_t)$. EMA sinyal periode 9 diinisialisasi dengan rata-rata sembilan nilai MACD pertama yang valid, dan histogram adalah selisih garis MACD dengan sinyal. Ketiga fitur MACD dibagi dengan $C_t$ agar skalanya relatif terhadap harga. Setelah awal data atau setiap celah yang tidak dapat dijelaskan, observasi baru digunakan setelah sedikitnya 50 bar perdagangan valid untuk pemanasan seluruh indikator.

\subsection{Fitur Makroekonomi}

Empat tipe CPI dipertahankan sebagai keluarga terpisah, demikian pula NFP dan tingkat pengangguran. Untuk setiap keluarga berwujud numerik tersebut, nilai kejutan rilis dihitung sebagai
\begin{equation}
u_{e}=A_{e}-F_{e},
\label{eq:kejutan-makro}
\end{equation}
dengan $A_e$ sebagai kolom \emph{actual} dan $F_e$ sebagai kolom \emph{forecast} kejadian $e$ pada snapshot MT5 beku. Kejutan tidak dihitung untuk keputusan suku bunga dan peristiwa FOMC karena kolom prakiraan pada arsip tersebut tidak tersedia. Kedua tipe peristiwa tersebut hanya menyumbangkan penanda kejadian pascarilis tanpa membentuk angka kejutan rekaan.

Untuk tiap keluarga disusun nilai kejutan terbaru, umur rilis, penanda kedaluwarsa, serta indikator bahwa keputusan berada dalam 1, 6, dan 24 jam setelah rilis. Umur dibatasi pada 168 jam untuk kebutuhan representasi; pada umur $\geq168$ jam, nilai kejutan menjadi hilang dan penanda kedaluwarsa diaktifkan. Angka dari satuan berbeda tetap berada dalam kolom terpisah dan dinormalisasi hanya berdasarkan data latih. Jika beberapa pengumuman memiliki waktu sama, urutan pemrosesan ditentukan secara tetap berdasarkan keluarga dan pengenal kejadian, kemudian seluruh informasi yang sah tersedia sebelum $\tau_t$ direpresentasikan tanpa bergantung pada urutan berkas. Jadwal peristiwa masa depan tidak digunakan dalam konfigurasi utama; seluruh penanda kalender dibentuk setelah kejadian berlangsung.

\section{Konfigurasi Eksperimen}

Tiga konfigurasi pada Tabel \ref{tab:konfigurasi-fitur} diberikan kepada kedua model. K1, K2, dan K3 dilaksanakan pada setiap periode uji utama yang memenuhi persyaratan evaluasi tiga kelas. Masker observasi akhir dibentuk sekali sebelum pelatihan berdasarkan aturan data, waktu, label, dan pemanasan. Untuk menjamin perbandingan yang setara, bar uji yang digunakan harus sama untuk seluruh konfigurasi dan kedua model; K1 atau K2 tidak memperoleh bar tambahan hanya karena tidak memerlukan fitur makro.

\begin{table}[htbp]
    \centering
    \caption{Konfigurasi fitur eksperimen}
    \label{tab:konfigurasi-fitur}
    \begin{tabular}{cll}
        \hline
        Kode & Kelompok fitur & Isi utama \\
        \hline
        K1 & Harga & 9 fitur harga \\
        K2 & Harga + teknikal & K1 + RSI dan 3 fitur MACD \\
        K3 & Harga + teknikal + makro & K2 + fitur snapshot kalender MT5 \\
        \hline
    \end{tabular}
\end{table}

Prediktor konstan yang selalu memilih kelas terbanyak pada data latih dihitung sebagai diagnostik kewajaran, bukan sebagai model ketiga. Jika jumlah kelas sama, pemilihan mengikuti urutan tetap BUY, HOLD, lalu SELL. Diagnostik ini juga menggunakan bar uji yang sama.

\section{Validasi Kronologis dan Penyetelan Model}

\subsection{Validasi Luar dan Dalam}

Evaluasi utama menggunakan \emph{expanding-window walk-forward}. Interval dinyatakan setengah terbuka dalam UTC, yaitu batas awal termasuk dan batas akhir tidak termasuk. Data latih awal adalah $[1$ Januari 2020, $1$ Januari 2023$)$. Tujuh periode uji utama dicantumkan pada Tabel \ref{tab:periode-utama}. Penyebutan semester merupakan nama blok kalender; blok terakhir secara khusus mencakup Januari--Juli 2026. Setelah suatu periode selesai diuji, observasi periode tersebut dapat masuk ke data latih untuk periode berikutnya, sesuai penggunaan model secara kronologis.

\begin{table}[htbp]
    \centering
    \caption{Periode uji kronologis data utama}
    \label{tab:periode-utama}
    \begin{tabular}{cll}
        \hline
        Periode uji & Awal (termasuk) & Akhir (tidak termasuk) \\
        \hline
        1 & 1 Januari 2023 & 1 Juli 2023 \\
        2 & 1 Juli 2023 & 1 Januari 2024 \\
        3 & 1 Januari 2024 & 1 Juli 2024 \\
        4 & 1 Juli 2024 & 1 Januari 2025 \\
        5 & 1 Januari 2025 & 1 Juli 2025 \\
        6 & 1 Juli 2025 & 1 Januari 2026 \\
        7 & 1 Januari 2026 & 1 Agustus 2026 \\
        \hline
    \end{tabular}
\end{table}

Pada setiap pelatihan luar, validasi internal menggunakan 18 bulan terakhir dari data latih sebagai tiga blok validasi enam bulanan. Seluruh riwayat sebelum awal tiap blok menjadi data latih internal yang mengembang. Untuk periode uji utama pertama, blok validasinya adalah Juli--Desember 2021, Januari--Juni 2022, dan Juli--Desember 2022. Pada periode uji berikutnya, jendela 18 bulan bergerak maju mengikuti akhir data latih luar.

\subsection{Purge dan Prapemrosesan pada Setiap Pemisahan}

Sebelum setiap blok validasi atau uji, seluruh observasi latih dengan akhir label $\tau_{t+1}$ yang sama dengan atau melewati waktu keputusan pertama pada blok berikutnya dibuang. Aturan ini menghasilkan \emph{purge} sedikitnya satu bar H1 dan diterapkan pada pemisahan internal maupun luar. Pemanasan fitur dapat membaca bar lampau sebelum awal blok karena bar tersebut telah tersedia; embargo sepanjang jendela fitur tidak ditambahkan. Setiap label uji tetap wajib berakhir di dalam interval uji.

Seluruh operasi yang mempelajari parameter, yaitu imputasi median, indikator kolom hilang, pembuangan kolom yang seluruhnya hilang, penskalaan, normalisasi satuan, bobot kelas, dan pemilihan hiperparameter, diestimasi ulang hanya dari data latih pada setiap blok validasi internal dan periode uji utama. Data latih harus memuat ketiga kelas. Jika salah satu kelas tidak ada, blok atau periode tersebut dinyatakan tidak memenuhi persyaratan evaluasi tiga kelas dan tidak diubah menjadi masalah dua kelas. Kondisi tersebut dicatat sebagai kegagalan evaluasi.

\subsection{Model dan Kisi Hiperparameter}

Regresi Logistik Multinomial menggunakan penalti L2, \emph{solver} \texttt{lbfgs}, dan maksimum 2.000 iterasi. Kandidat $C$ adalah $\{0{,}01;0{,}1;1;10\}$. Fitur distandardisasi dari data latih. Bobot setiap kelas $c$ dihitung sebagai
\begin{equation}
w_c=\frac{n}{3n_c},
\label{eq:bobot-kelas}
\end{equation}
dengan $n$ jumlah observasi latih dan $n_c$ jumlah observasi kelas $c$ pada data latih.

XGBoost menggunakan objektif \texttt{multi:softprob} untuk tiga kelas. Delapan kandidat pada kisi pencarian merupakan hasil kali Kartesius parameter pada Tabel \ref{tab:kisi-xgboost}. Nilai \texttt{n\_estimators} menyatakan jumlah iterasi \emph{boosting}, bukan jumlah literal seluruh pohon pada model multikelas.

\begin{table}[htbp]
    \centering
    \caption{Kisi hiperparameter XGBoost}
    \label{tab:kisi-xgboost}
    \begin{tabular}{lll}
        \hline
        Parameter & Kandidat & Keterangan \\
        \hline
        \texttt{max\_depth} & $\{3,5\}$ & Kedalaman maksimum \\
        \texttt{learning\_rate} & $\{0{,}05;0{,}1\}$ & Laju pembelajaran \\
        \texttt{n\_estimators} & $\{100,300\}$ & Iterasi \emph{boosting} \\
        \hline
    \end{tabular}
\end{table}

Setiap kandidat XGBoost menggunakan proporsi pengambilan observasi (\texttt{subsample}) sebesar 0,8 dan proporsi fitur per pohon (\texttt{colsample\_bytree}) sebesar 0,8. Regularisasi L2 ditetapkan pada $\lambda=1$ dengan benih acak 42. Bobot observasi mengikuti Persamaan \ref{eq:bobot-kelas}. \emph{Early stopping} tidak digunakan. Kebijakan nilai hilang yang sama diterapkan melalui alur prapemrosesan Regresi Logistik dan XGBoost; kemampuan bawaan XGBoost untuk menerima nilai hilang tidak digunakan untuk mengubah definisi sampel.

\subsection{Penyetelan Hiperparameter dan Analisis Sensitivitas}

Kisi yang sama digunakan untuk setiap konfigurasi dalam model yang sama. Kandidat dipilih berdasarkan rata-rata macro-F1 pada tiga blok internal, kemudian model dipasang ulang pada seluruh data latih luar. Jika skor sama, Regresi Logistik memilih $C$ yang lebih kecil. XGBoost berturut-turut memilih kedalaman lebih kecil, jumlah iterasi \emph{boosting} lebih sedikit, lalu \emph{learning rate} lebih kecil. Kandidat yang gagal atau tidak konvergen pada salah satu blok validasi internal tidak memenuhi syarat pemilihan; rata-rata skor tidak dihitung hanya dari blok yang tersisa. Jika seluruh kandidat gagal, model pada periode uji tersebut dicatat sebagai kegagalan evaluasi beserta penyebabnya. Seluruh kegagalan tetap dicatat.

Perbandingan utama membolehkan hiperparameter berbeda antarkonfigurasi sehingga hasilnya ditafsirkan sebagai fitur ditambah penyetelan ulang. Dalam analisis sensitivitas parameter tetap, hiperparameter K1 yang dipilih pada validasi internal digunakan kembali untuk K2 dan K3 pada model dan periode uji utama yang sama. Analisis ini dilaporkan terpisah dan tidak menggantikan analisis utama.

\section{Metrik dan Analisis Hasil}

Macro-F1 menjadi metrik utama dan balanced accuracy menjadi metrik pendamping sebagaimana didefinisikan pada Bab 2. Matriks kebingungan selalu menggunakan urutan tiga label BUY, HOLD, dan SELL. Jika penyebut F1 suatu kelas nol, F1 kelas tersebut didefinisikan nol sehingga penyebut rata-rata macro-F1 tetap tiga kelas. Jika satu kelas benar tidak muncul pada data uji, recall kelas tersebut tidak terdefinisi; akibatnya balanced accuracy pada periode uji tersebut dinyatakan tidak terdefinisi dan tidak disajikan sebagai balanced accuracy tiga kelas yang sebanding. Kondisi ini, beserta jumlah kelas dan jumlah observasi, dilaporkan secara eksplisit.

Untuk setiap model, metrik, dan periode uji yang sah, perubahan skor dihitung dengan Persamaan \ref{eq:delta_teknikal} dan \ref{eq:delta_makro}. Ringkasan perubahan berpasangan meliputi rata-rata tanpa pembobotan jumlah bar, median, simpangan baku sampel, serta jumlah periode uji dengan perubahan positif. Ringkasan dihitung hanya untuk pasangan skor yang terdefinisi dan jumlah pasangan turut dilaporkan; simpangan baku tidak dihitung jika tersedia kurang dari dua pasangan. Penyebut jumlah periode uji selalu ditampilkan, termasuk periode uji yang tidak memenuhi persyaratan evaluasi, balanced accuracy yang tidak terdefinisi, dan kegagalan kandidat. Periode uji yang gagal tetap dicantumkan dalam pelaporan. Karena hanya terdapat tujuh periode uji utama yang saling bergantung akibat jendela latih mengembang, analisis bersifat deskriptif dan tidak menggunakan uji hipotesis atau nilai-$p$ naif yang mengasumsikan independensi periode uji.

Perubahan positif, nol, maupun negatif sama-sama menjawab pertanyaan penelitian. Konsistensi dinilai berdasarkan arah perubahan pada periode uji yang tersedia dan pada kedua model, tanpa mengubah kriteria kelayakan berdasarkan hasil skor. Hasil tidak ditafsirkan sebagai keuntungan transaksi, dampak kausal ekonomi, atau kontribusi individual setiap fitur.

\section{Reproduksibilitas dan Ancaman Validitas}

Setiap eksekusi menyimpan hash data masukan dan keluaran, versi perangkat lunak, benih acak, batas waktu setiap periode uji, aturan eksklusi, daftar kolom yang dibuang, parameter terpilih, status konvergensi, prediksi per bar, label benar, dan metrik. Data mentah, data turunan, dan artefak hasil ditempatkan secara terpisah sehingga proses eksperimen tidak menimpa data sumber. Catatan tersebut memungkinkan hasil ditelusuri tanpa menganggap bahwa kesamaan kode menjamin kesamaan data vendor.

Ancaman validitas utama berasal dari ketiadaan snapshot \emph{point-in-time} independen yang merekam perubahan nilai kalender. Nilai \emph{actual} dan \emph{forecast} historis yang diperoleh secara retrospektif dari MT5 mungkin mencerminkan koreksi setelah publikasi. Oleh sebab itu, hasil K3 berlaku untuk snapshot kalender beku yang digunakan dan tidak disetarakan dengan simulasi waktu nyata yang sepenuhnya \emph{point-in-time}. Ancaman lain meliputi jarangnya peristiwa makroekonomi, perubahan rezim harga, karakteristik data pada server MetaQuotes-Demo, serta ketergantungan antarperiode uji. Kesimpulan dibatasi pada XAUUSD H1 dari MetaQuotes-Demo, rentang waktu, definisi label, dua model, dan konfigurasi fitur yang dinyatakan pada penelitian ini.

\section{Jadwal Penelitian}

Jadwal penelitian disusun agar Penilaian Kemajuan 2 (P2) selesai sebelum Desember 2026. Waktu setelah P2 dialokasikan untuk menindaklanjuti masukan Pembimbing, menyelesaikan naskah dan artikel ilmiah, memenuhi persyaratan administratif, serta mempersiapkan Seminar Hasil. Eksperimen K3 dan analisis $K3-K2$ ditetapkan sebagai luaran penelitian. Rencana pelaksanaan penelitian ditunjukkan pada Tabel \ref{tab:jadwal-penelitian}.

\begin{table}[htbp]
    \centering
    \caption{Jadwal pelaksanaan penelitian}
    \label{tab:jadwal-penelitian}
    \small
    \setlength{\tabcolsep}{3pt}
    \begin{tabular}{>{\raggedright\arraybackslash}p{4.0cm}>{\raggedright\arraybackslash}p{3.0cm}>{\raggedright\arraybackslash}p{6.1cm}}
        \hline
        Kegiatan & Periode & Target \\
        \hline
        P0 dan pembekuan protokol & 21 September 2026 & Protokol serta target P1--P2 disepakati \\
        Implementasi pipeline & 22 September--4 Oktober 2026 & Pipeline validasi kronologis teruji \\
        Eksperimen K1--K2 & 5--18 Oktober 2026 & Prediksi, metrik, dan analisis sensitivitas \\
        Analisis awal dan P1 & 19--25 Oktober 2026 & Capaian sedikitnya 50\% dan ringkasan $K2-K1$ \\
        Eksperimen K3 & 26 Oktober--8 November 2026 & Prediksi, metrik, dan analisis $K3-K2$ \\
        Analisis dan penulisan & 9--22 November 2026 & Hasil, pembahasan, dan draf naskah lengkap \\
        Finalisasi dan P2 & 23--27 November 2026 & Capaian sedikitnya 80\% dan berita acara P2 \\
        Persiapan Seminar Hasil & 28 November--11 Desember 2026 & Naskah final, artikel, berkas, dan presentasi \\
        \hline
    \end{tabular}
\end{table}
